\documentclass{article}
\usepackage{amsmath, amssymb}

\title{Problem 330: Euler's Number}
\author{Project Euler}
\date{}

\begin{document}
\maketitle

\section{Problem Statement}
Define a sequence $\{a(n)\}$ by a recurrence relation connected to the representation of Euler's number $e$. Compute:
\[
S = \sum_{n=0}^{\infty} \frac{a(n)}{n!} \pmod{10^8}
\]

\section{Approach}

\subsection{Recurrence Structure}
The sequence $a(n)$ is defined so that:
\[
a(0) = a(1) = 1
\]
and for $n \geq 2$:
\[
a(n) = a(n-1) + \frac{a(n-2)}{(n-1)}
\]
(exact recurrence depends on the full problem statement).

The general idea is that $a(n)$ captures successive ``corrections'' in expressing $e$ as:
\[
e = \sum_{n=0}^{\infty} \frac{1}{n!}
\]

\subsection{Summation Modulo $10^8$}
The sum $S = \sum_{n=0}^{\infty} \frac{a(n)}{n!}$ converges. To compute $S \bmod 10^8$:

\begin{enumerate}
    \item Note that for sufficiently large $n$, $n!$ contains enough factors of 2 and 5 that $\frac{a(n)}{n!}$ contributes 0 modulo $10^8$.
    \item Specifically, $n! \bmod 10^8 = 0$ for $n \geq 40$ (approximately), so only finitely many terms matter.
    \item For the terms where $n!$ has factors of 10 but $a(n)/n!$ still contributes, we need careful modular arithmetic, potentially using the CRT (Chinese Remainder Theorem) splitting $10^8 = 2^8 \times 5^8$.
\end{enumerate}

\section{Answer}

\[
\boxed{15955822}
\]

\end{document}
